\documentclass[12pt,a4paper]{article}
\usepackage[utf8]{inputenc}
\usepackage[T1]{fontenc}
\usepackage{amsmath}
\usepackage{amssymb}
\usepackage{amsfonts}
\usepackage{mathtools}
\usepackage{physics}
\usepackage{siunitx}
\usepackage{graphicx}
\usepackage{xcolor}
\usepackage{hyperref}
\usepackage{geometry}
\usepackage{microtype}
\usepackage{lmodern}
\geometry{margin=2.5cm}
\hypersetup{colorlinks=true,linkcolor=blue,urlcolor=blue}

\title{WhatsApp Image 2026-03-26 at 21.55.23.jpeg}
\author{Biko}
\date{03 mai 2026}

\begin{document}
\maketitle

\section*{WhatsApp Image 2026-03-26 at 21.55.23.jpeg}

PROBLEM IN ONE DIMENSION CHAP. 5

5-10 Hermite polynomials. The $n$th Hermite polynomial $H_{n}(x)$ is a solution of the differential equation
$$H_{n}^{\prime \prime}-2 x H_{n}^{\prime}+2 n H_{n}=0$$
it is a polynomial of degree $n$ in which the coefficient of the term in $x^{n}$ is $2^{n}$. By differentiation of Eq. (5-116), we obtain
$$\left(H_{n}^{\prime}\right)^{\prime \prime}-2 x\left(H_{n}^{\prime}\right)^{\prime}-2 H_{n}^{\prime}+2 n H_{n}^{\prime}=0$$
whence the function $H_{n}^{\prime}(x)=d H_{n} / d x$ satisfies the differential equation
$$\phi^{\prime \prime}-2 x \phi^{\prime}+2(n-1) \phi^{\prime}=0$$

This equation, however, is also satisfied by $H_{n-1}$, and since it has only one polynomial solution, we conclude that
$$H_{n}^{\prime}=C H_{n-1} .$$

By comparison of the terms in $x^{n-1}, C$ can be evaluated:
$$n 2^{n} x^{n-1}=C 2^{n-1} x^{n-1}$$

The constant $C$ is therefore $2 n$, and
$$H_{n}^{\prime}=2 n H_{n-1},$$
which is a recurrence formula. By further differentiation and manipulation of Eqs. (5-120) and (5-116), the relations
$$\begin{array}{c}
H_{n+1}-2 x H_{n}+2 n H_{n-1}=0, \\
H_{n+1}=2 x H_{n}-H_{n}^{\prime}
\end{array}$$
can be derived (Problem 5-27). Equation (5-122) can be used to construct the polynomials in consecutive order, beginning with $H_{0}=1$.

The generating function, defined by
$$g(x, h)=\sum_{n=0}^{\infty} H_{n}(x) \frac{h^{n}}{n!},$$
is of great usefulness in calculations that involve the Hermite functions. By means of the recurrence formulae for $H_{n}$, a closed expression for $g(x, h)$ can be derived. Differentiating partially with respect to $h$, one obtains


\end{document}
